% !TEX root = ../main.tex
\chapter{米田の補題：インターフェースから値を見る}
\label{chap:yoneda-interface}

\begin{chapterabstract}
本章では、米田の補題を本格的な定理証明としてではなく、「値はそれに対するすべての観測方法から特徴づけられる」という直感として導入する。対象の内部表現を直接のぞくのではなく、その対象をどのように使えるか、どのような関数で観測できるかを見る。この見方は、API設計、抽象データ型、テストダブル、観測可能性、インターフェースの安定性を考えるうえで役に立つ。

圏論側では、集合の圏における小さな影として
\[
  \operatorname{Nat}(\operatorname{Hom}(A,-), \mathsf{Id}) \cong A
\]
を読む。Lean側では、完全な Mathlib の圏論形式化には進まず、値を関数で観測する小さな例を扱う。
\end{chapterabstract}

\begin{learninggoals}
\item 米田の補題を「観測方法と値の対応」という直感で説明できる。
\item 代表可能関手と自然変換が、ソフトウェア上のインターフェースにどう似ているかを説明できる。
\item Leanで「すべての観測が一致するなら値が一致する」という小さな影を表せる。
\item API設計、テストダブル、観測可能性に対して、米田的な見方を安全な比喩として使える。
\end{learninggoals}

\section{この章を読む理由}

実務では、値の内部表現を直接見られないことが多い。ライブラリの利用者は、公開されたメソッドだけを使う。APIの利用者は、レスポンスの公開フィールドだけを見る。テストコードは、本物のデータベースではなく、インターフェース越しの振る舞いを見る。つまり、多くの場面で私たちは「値そのもの」ではなく、「値をどう観測できるか」を扱っている。

\begin{intuitionbox}
米田の補題を本章の言葉で読むと、次のようになる。

あるものが何であるかは、そのものをあらゆる方法で使ったときにどう振る舞うかによって決まる。内部表現を開けて見る代わりに、すべての観測窓から見える振る舞いを集めると、そのものを復元できる。
\end{intuitionbox}

この言い方は、まだ厳密な米田の補題ではない。しかし、ソフトウェアエンジニアが最初に持つべき直感としては十分に有用である。とくに、隠蔽された実装、公開インターフェース、adapter、テストダブル、観測ログを考えるときに、この見方は「何を同じと見なすのか」を整理してくれる。

\FigurePlaceholder{fig-ch38-yoneda-overview.png}{0.90}{値を直接見るのではなく、すべての観測方法から値を特徴づける}{fig:ch38-yoneda-overview}

図\ref{fig:ch38-yoneda-overview}では、中央にある値を直接のぞくのではなく、複数の観測関数を通じて外から見る状況を描く。観測方法が十分に豊富であれば、外からの振る舞いだけで値を区別できる。

\section{本章では直感だけを扱う}

米田の補題は、圏論の中心的な定理である。正確に扱うには、圏、関手、自然変換、代表可能関手、自然同型を組み合わせる必要がある。本書ではすでに関手と自然変換を扱ったが、本章で完全な証明までは行わない。

\begin{warningbox}
本章の目的は、米田の補題を使って高度な圏論計算をすることではない。目的は、「対象を、その対象に関するすべての使い方から見る」という発想を得ることである。

Leanでも、Mathlibの \texttt{CategoryTheory} を使った本格的な米田埋め込みや米田補題の形式化には進まない。それは第40章以降の発展課題である。
\end{warningbox}

この境界を明確にしておくことは重要である。比喩は理解を助けるが、比喩だけで定理を証明したことにはならない。本章では、比喩と最小モデルを足場にして、後で標準的な圏論教材へ進める入口を作る。

\section{観測としての関数}

まず、集合と関数の世界で考える。型 \texttt{A} の値 \texttt{x : A} があるとする。この値を観測するとは、何らかの関数
\[
  f : A \to B
\]
を適用して、結果 \texttt{f x : B} を見ることだと考える。

たとえば、\texttt{User} 型の値があるとき、次のような観測が考えられる。

\begin{center}
\begin{tabular}{lll}
\toprule
観測関数 & 返すもの & 実務上の意味 \\
\midrule
\texttt{userId : User -> Nat} & ID & 同一性確認に使う \\
\texttt{displayName : User -> String} & 表示名 & UIに出す \\
\texttt{isActive : User -> Bool} & 有効状態 & 権限制御に使う \\
\texttt{toJson : User -> String} & JSON表現 & API境界で使う \\
\bottomrule
\end{tabular}
\end{center}

\begin{softwarebox}
インターフェースは、値の内部を直接公開するものではない。むしろ、「利用者がその値に対して何をできるか」を並べたものである。米田的な見方では、この「できることの全体」が値の意味を決める。
\end{softwarebox}

ここで重要なのは、観測が一つだけでは足りない場合があることだ。IDだけが同じユーザーは、名前や権限が違うかもしれない。名前だけが同じユーザーは、別人かもしれない。観測が少なすぎると、外から区別できない値が残る。

\section{すべての観測方法が値を決める}

観測方法を一つだけでなく、すべて考えたらどうなるだろうか。型 \texttt{A} の値 \texttt{x} に対して、任意の型 \texttt{B} と任意の関数 \texttt{f : A -> B} を選び、\texttt{f x} を見る。

もし二つの値 \texttt{x y : A} について、すべての \texttt{B} とすべての \texttt{f : A -> B} で
\[
  f(x) = f(y)
\]
が成り立つなら、\texttt{x} と \texttt{y} は同じ値である。なぜなら、観測関数として恒等関数 \texttt{id : A -> A} も選べるからである。恒等関数で観測すれば、\texttt{id x = id y}、つまり \texttt{x = y} が得られる。

\begin{definitionbox}
本章でいう「すべての観測方法」とは、値 \texttt{x : A} に対して適用できるすべての関数 \texttt{f : A -> B} を考える、という意味である。

この考えは、集合の圏 \(\mathsf{Set}\) における米田補題の非常に小さな影である。一般の圏では、\texttt{A -> B} の代わりに射 \(A \to B\) を考える。
\end{definitionbox}

この事実だけを見ると、恒等関数を使った当たり前の話に見える。しかし米田の補題の力は、「この当たり前」を一般の圏、一般の関手、自然変換の言葉で成り立たせるところにある。

\section{米田の補題を一文で読む}

米田の補題の標準形とvarianceは、Hom関手からの自然変換と対象上の要素の自然な対応として扱う~\cite{riehl-category-theory-in-context}。本章後半のobserver例は、その定理そのものではなく値分離の小さな影である。

圏 \(\mathcal{C}\) の対象 \(A\) を固定する。各対象 \(X\) に対して、\(A\) から \(X\) への射全体 \(\mathcal{C}(A,X)\) を対応させる。これは、\(A\) から外へ向かう使い方を集める関手である。

この関手を、\(\mathcal{C}(A,-)\) と書く。\(-\) は「ここに任意の対象を入れる」という意味である。

\begin{definitionbox}
\textbf{代表可能関手}とは、ある対象 \(A\) からの射全体、またはある対象 \(A\) への射全体によって作られる関手である。本章では、\(A\) から各対象への射を集める \(\mathcal{C}(A,-)\) を使う。
\end{definitionbox}

米田の補題の共変版は、ざっくり次の形で読める。
\[
  \operatorname{Nat}(\mathcal{C}(A,-), F) \cong F(A)
\]
左辺は、\(A\) からのすべての使い方を、関手 \(F\) が与えるデータへ一貫して変換する方法の集まりである。右辺は、\(F\) を \(A\) に適用して得られるデータそのものである。

\begin{intuitionbox}
米田の補題は、「\(A\) から始まるすべての使い方を一貫して解釈する方法」は、「\(A\) における値」を一つ選ぶことと同じだ、と言っている。
\end{intuitionbox}

ここで「一貫して」が重要である。対象 \(X\) ごとにばらばらな対応を選ぶのではない。射をさらに合成したときにも、対応が合成と干渉しない必要がある。この一貫性が自然性である。

\FigurePlaceholder{fig-ch38-natural-observation.png}{0.92}{観測してから変換する経路と、変換してから観測する経路が一致する自然性}{fig:ch38-natural-observation}

図\ref{fig:ch38-natural-observation}では、観測関数 \(f : A \to B\) と、その後処理 \(g : B \to C\) を考える。先に \(f\) で観測してから \(g\) を適用することと、最初から \(g \circ f\) で観測することは同じである。この当たり前の可換性が、自然性の最小の雰囲気である。

\section{集合の圏での小さな影}

\(\mathsf{Set}\) で \(F\) を恒等関手 \(\mathsf{Id}\) とすると、米田の形は次のようになる。
\[
  \operatorname{Nat}(\operatorname{Hom}(A,-), \mathsf{Id}) \cong A
\]
これは、\(A\) の値 \(x\) と、「任意の関数 \(f : A \to B\) に対して \(f(x)\) を返す仕組み」が対応する、という意味で読める。

値 \(x : A\) があれば、任意の \(B\) と \(f : A \to B\) に対して、\(f(x)\) を返せる。逆に、そのような仕組みがあれば、\(B = A\)、\(f = \operatorname{id}\) を選ぶことで、元の値を取り出せる。

\begin{softwarebox}
これは、抽象データ型のインターフェースに似ている。値そのものを直接受け取らなくても、「どんな関数で観測しても、その結果を返してくれるもの」があれば、少なくとも十分豊かな世界では、その値を持っているのと同じ情報を持っている。
\end{softwarebox}

ただし、この話は「すべての観測方法」を考える点に注意する。現実のAPIでは、公開された観測方法は有限であり、権限やパフォーマンスや互換性の都合で制限される。そのため、現実のインターフェースだけで常に内部値が完全に決まるとは限らない。

\section{Lean 4で観測の影を見る}

Leanでは、値 \texttt{x : A} から「任意の型 \texttt{B} と関数 \texttt{A -> B} に対して \texttt{B} の値を返すもの」を定義できる。これは、値を観測関数に渡すだけの小さな関数である。

\begin{leanbox}
次のコードでは、\texttt{observeAll x} が「任意の観測関数 \texttt{f} に対して \texttt{f x} を返すもの」を表す。\texttt{same\_by\_all\_observers} は、すべての観測が一致するなら値が一致する、という小さな事実を示す。
\end{leanbox}

\begin{listing}[htbp]
\begin{minted}{lean4}
def observeAll {A : Type} (x : A) :
    (B : Type) -> (A -> B) -> B :=
  fun _ f => f x

theorem recover_by_id {A : Type} (x : A) :
    observeAll x A (fun a => a) = x :=
  rfl

theorem same_by_all_observers {A : Type} {x y : A}
    (h : forall (B : Type) (f : A -> B), f x = f y) :
    x = y :=
  h A (fun a => a)
\end{minted}
\caption{すべての観測が一致するなら値が一致する}
\label{lst:ch38_all_observers}
\end{listing}

ここで \path{recover_by_id} は、恒等関数で観測すれば元の値が返ることを示している。\path{same_by_all_observers} は、すべての観測関数で結果が同じなら、恒等関数で見ても同じなので、値そのものが等しいことを示している。

次に、観測してから後処理することと、後処理込みの観測を一度に行うことが一致する様子を見る。

\begin{listing}[htbp]
\begin{minted}{lean4}
def post {A B C : Type} (g : B -> C) (f : A -> B) :
    A -> C :=
  fun a => g (f a)

theorem equal_values_same_observations
    {A _B : Type} {x y : A}
    (h : x = y) (f : A -> B) :
    f x = f y := by
  cases h
  rfl

theorem observation_commutes {A B C : Type} (x : A)
    (f : A -> B) (g : B -> C) :
    g (observeAll x B f) =
      observeAll x C (post g f) :=
  rfl
\end{minted}
\caption{観測と後処理の順序が一致する}
\label{lst:ch38_observation_commutes}
\end{listing}

このコードは、一般圏論としての自然性を証明しているわけではない。しかし、自然性の雰囲気は見える。先に \texttt{f} で観測してから \texttt{g} で後処理することと、\texttt{g} と \texttt{f} を合成した観測を直接使うことは同じである。

\section{部分的なインターフェースでは足りない}

現実のAPIやクラスでは、すべての関数 \texttt{A -> B} を公開することはない。公開するのは、ごく一部のメソッドやフィールドである。この場合、「公開インターフェースから区別できない」ことと「内部値が同じである」ことは一致しない。

\FigurePlaceholder{fig-ch38-partial-interface.png}{0.88}{部分的なインターフェースでは区別できない値が残る}{fig:ch38-partial-interface}

図\ref{fig:ch38-partial-interface}では、二つの値が同じIDを返すが、名前は異なる状況を描く。IDだけを公開するインターフェースから見ると二つの値は同じに見える。しかし、名前も観測できるなら区別できる。

\begin{listing}[htbp]
\begin{minted}{lean4}
structure User where
  id : Nat
  name : String

structure UserView where
  id : Nat
  name : String

def toView (u : User) : UserView :=
  { id := u.id, name := u.name }

theorem user_eq_from_view_eq {u v : User}
    (h : toView u = toView v) : u = v := by
  cases u
  cases v
  cases h
  rfl

def byId (u : User) : Nat := u.id

def sameById (u v : User) : Prop :=
  byId u = byId v

def userA : User := { id := 1, name := "Alice" }
def userB : User := { id := 1, name := "Alicia" }

example : sameById userA userB :=
  rfl
\end{minted}
\caption{IDだけの観測ではユーザー全体は決まらない}
\label{lst:ch38_partial_observer}
\end{listing}

\texttt{toView} は、\texttt{id} と \texttt{name} の両方を公開するビューである。このビューが等しいなら、今回の単純な \texttt{User} では値そのものも等しい。一方、\texttt{byId} だけを見る \texttt{sameById} では、\texttt{userA} と \texttt{userB} を区別できない。

\begin{warningbox}
「インターフェースから見て同じ」と言うときは、どのインターフェースから見るのかを明示する必要がある。観測方法が少なければ、異なる値が同じに見えることがある。これは抽象化として有用な場合もあるが、完全な同一性の証明ではない。
\end{warningbox}

\section{API設計への読み替え}

米田的な見方は、API設計で次の問いを立てる助けになる。

\begin{itemize}
\item 利用者は、この値をどの観測方法で見るのか。
\item その観測方法は、区別すべき値を本当に区別できるのか。
\item 逆に、区別しなくてよい内部差分を隠せているのか。
\item adapterを通しても、観測結果は一貫しているのか。
\end{itemize}

たとえば、\texttt{User} の内部表現を変更しても、公開APIが返す \texttt{id}、\texttt{name}、\texttt{role} がすべて保存されるなら、利用者から見た互換性は保たれるかもしれない。これは「公開された観測方法に関して同値」という主張である。

\begin{softwarebox}
実務上は、完全な内部同一性よりも、公開インターフェースに関する観測同値が重要なことが多い。API互換性、SDK adapter、テストダブルは、内部表現ではなく、公開された使い方の一貫性で評価される。
\end{softwarebox}

この見方は、第31章の自然性にもつながる。adapterが業務ロジックと干渉しないとは、観測や変換の順序を入れ替えても結果が一致するということである。

\section{テストダブルと観測可能性}

テストダブルは、本物の実装と同じ内部構造を持つ必要はない。必要なのは、テストが観測する範囲で同じ振る舞いをすることである。これは、観測可能な振る舞いに関する同値である。

たとえば、本物のメール送信サービスは外部SMTPサーバーに接続するかもしれない。しかしテストダブルは、送信された宛先と件名をメモリに記録するだけでよい。テストが見る観測方法が「何が送信されたか」だけなら、その観測に関して同じ振る舞いをすれば十分である。

\begin{intuitionbox}
米田的な視点では、テストダブルの正しさは「内部が本物と同じか」ではなく、「必要な観測方法から見て本物と区別できないか」として考えられる。
\end{intuitionbox}

観測可能性も同じである。ログ、メトリクス、トレースは、実行中のシステムを外から見る観測関数のように振る舞う。観測方法が不十分なら、障害の原因を区別できない。観測方法が適切なら、内部状態を直接見なくても、問題を推定できる。

\section{表現可能性への橋渡し}

本章では、代表可能関手を「ある対象からの使い方を全部集めるもの」として紹介した。これは、圏論の重要語である表現可能性への入口である。

ある関手 \(F\) が代表可能であるとは、ある対象 \(A\) が存在して、\(F\) が \(\mathcal{C}(A,-)\) と同じ情報を持つ、ということである。ソフトウェア的には、「ある仕様や振る舞い全体が、一つの代表的なデータ構造からの使い方として説明できる」と読める。

\begin{definitionbox}
\textbf{表現可能性}とは、抽象的に与えられた振る舞いや問い合わせの集まりが、ある具体的な対象からの射の集まりとして表せることをいう。本章では詳細には踏み込まないが、米田の補題はこの表現可能性を扱うための基本道具になる。
\end{definitionbox}

表現可能性を本格的に扱うと、普遍性、随伴、極限、データベース的な問い合わせ、型理論的な意味論へつながる。ここから先は、標準的な圏論教材や Mathlib の \texttt{CategoryTheory} を読む段階で再び出会う。

\section{よくある誤解}

\begin{warningbox}
\textbf{誤解1：米田の補題は「ただのインターフェース設計の話」である。}

そうではない。インターフェース設計は、本章での比喩である。米田の補題そのものは、圏、関手、自然変換に関する厳密な定理である。

\medskip
\textbf{誤解2：一つの観測関数で同じなら値も同じである。}

一般には成り立たない。IDだけ、名前だけ、ログ出力だけのような部分的な観測では、異なる値を区別できないことがある。

\medskip
\textbf{誤解3：内部表現を隠せば常に安全である。}

隠蔽は有用だが、公開された観測方法が仕様上必要な違いを区別できなければ、バグや互換性問題を見逃すことがある。

\medskip
\textbf{誤解4：Leanで小例を書けば米田の補題を証明したことになる。}

本章のLeanコードは集合の圏における直感の影である。一般の圏論としての米田補題は、Mathlibの圏論APIや標準教材を使って別途学ぶ必要がある。
\end{warningbox}

\section{発展への道}

この章の先には、二つの学習ルートがある。

一つ目は、圏論側のルートである。代表可能関手、自然同型、米田埋め込み、前層、Kan拡張へ進む。米田の補題は、これらの話題の中心にある。

二つ目は、Lean側のルートである。Mathlib の \texttt{CategoryTheory} では、圏、関手、自然変換、同型が一般的な形で定義されている。本章の自作ミニモデルを、そのAPIへ移すには、universe、型クラス、記法、既存定理の検索に慣れる必要がある。

本書では、第40章で Mathlib の \texttt{CategoryTheory} へ進むための地図を示す。本章の目的は、その前に「なぜそんな抽象的な定理がインターフェースや観測の話と関係するのか」を直感的につかむことである。

\section{本章のまとめ}

米田の補題は、本章では「値や対象は、それに対するすべての観測方法から特徴づけられる」という直感として読んだ。

集合の圏では、値 \texttt{x : A} から「任意の \texttt{f : A -> B} に対して \texttt{f x} を返すもの」を作れる。

すべての観測が一致するなら、恒等関数で観測することで値そのものの一致が得られる。

現実のAPIやインターフェースでは、観測方法は一部だけなので、観測同値と内部同一性を区別する必要がある。

テストダブル、観測可能性、API互換性は、内部表現ではなく公開された観測方法に関する同値として考えられる。

本章のLeanコードは米田の補題の完全証明ではなく、発展的なMathlib圏論へ進むための最小モデルである。
